\subsection{Summary}

The goal of this chapter was to understand how a classical optimization problem can be transformed into a quantum optimization problem.

\highspace
We started from the \important{Quadratic Unconstrained Binary Optimization\break (QUBO)} formulation (\autopageref{sec:qubo-problems}):
\begin{equation*}
    \min_{\mathbf{x}} \mathbf{x}Q\mathbf{x}^{T}
\end{equation*}
Where the unknown variables are binary:
\begin{equation*}
    x_i \in \{0,1\}
\end{equation*}
A QUBO problem describes an energy landscape in which each binary configuration has an associated cost (or energy). The objective is to find the configuration with minimum energy. QUBO is extremely important because many combinatorial optimization problems can be expressed in this form.

\highspace
Since QUBO problems are \textbf{unconstrained}, \textbf{constraints cannot be imposed directly}. Instead, they are incorporated through penalty terms, which increase the energy of infeasible solutions while leaving feasible solutions unchanged. In this way, the optimization process naturally favors valid configurations.

\highspace
The key step is establishing a connection between optimization and quantum mechanics. Each binary string (\autopageref{sec:encoding-bit-strings-into-quantum-states}):
\begin{equation*}
    \mathbf{x}\in\{0,1\}^n
\end{equation*}
Is mapped to a computational basis state:
\begin{equation*}
    \mathbf{x}
    \longrightarrow
    \ket{x}
\end{equation*}
Using this correspondence, the QUBO energy function can be encoded into a Hamiltonian (\autopageref{sec:hamiltonian-recap}) whose eigenvalues coincide with the energies of the original optimization problem. Consequently, minimizing the QUBO becomes equivalent to finding the ground state of the Hamiltonian.
\begin{equation*}
    \arg\min E(\mathbf{x})
    \quad \Longleftrightarrow \quad
    \text{ground state of } \hat{H}
\end{equation*}

\highspace
To better describe physical optimization systems, the problem is rewritten using the \definition{Ising model} (\autopageref{sec:ising-model}). Instead of binary variables:
\begin{equation*}
    x_i\in\{0,1\}
\end{equation*}
The Ising formulation uses spin variables:
\begin{equation*}
    s_i\in\{-1,+1\}
\end{equation*}
The classical Ising energy is:
\begin{equation*}
    H_{\text{Ising}}
    =
    \sum_i h_i s_i
    +
    \sum_{i>j} J_{ij}s_is_j
\end{equation*}
Where:
\begin{itemize}
    \item $h_i$ are the local fields (biases),
    \item $J_{ij}$ are the interaction strengths (couplers).
\end{itemize}
Thus, QUBO and Ising are simply two equivalent ways of representing the same optimization problem.

\highspace
The classical spins are then replaced by quantum operators, producing the quantum Ising Hamiltonian (\autopageref{eq:quantum-ising-hamiltonian}):
\begin{equation*}
    \hat H_{\text{Ising}}
    =
    \sum_i h_i Z_i
    +
    \sum_{i>j} J_{ij} Z_i Z_j
\end{equation*}
In this formulation, the eigenstates of the Hamiltonian correspond to candidate solutions, while the ground state corresponds to the optimal solution.

\highspace
Finding the ground state directly is generally computationally difficult. This motivates the use of \definition{Quantum Annealing} (\autopageref{sec:quantum-annealing}). The system is initially prepared in the ground state of a simple Hamiltonian:
\begin{equation*}
    H_A = -\sum_i X_i
\end{equation*}
Whose ground state is easy to construct. The Hamiltonian is then slowly transformed into the problem Hamiltonian that encodes the optimization problem.

\highspace
The complete annealing process is described by the \important{Transverse-Field Ising Hamiltonian}:
\begin{equation*}
    H(s)
    =
    -\frac{A(s)}{2}
    \left(\sum_i \sigma_x^{(i)}\right)
    +
    \frac{B(s)}{2}
    \left(
        \sum_i h_i \sigma_z^{(i)}
        +
        \sum_{i>j} J_{ij}\sigma_z^{(i)}\sigma_z^{(j)}
    \right)
\end{equation*}
Initially, the transverse-field term dominates and creates a superposition over many possible solutions. As the annealing progresses, the problem Hamiltonian gradually takes over until the final state encodes the solution of the optimization problem.

\highspace
The success of quantum annealing is explained through the evolution of the eigenspectrum (\autopageref{sec:evolution-of-eigenspectrum}). During the annealing process, the Hamiltonian possesses a sequence of energy levels:
\begin{equation*}
    E_0(s)
    \leq
    E_1(s)
    \leq
    E_2(s)
    \leq
    \cdots
\end{equation*}
Where $E_0(s)$ is the instantaneous ground-state energy. If the evolution is sufficiently slow and the energy gap between the ground state and the excited states remains large enough, the system remains in the ground state throughout the evolution and eventually reaches the optimal solution.

\highspace
In summary, the complete chain saw in this chapter is:
\begin{enumerate}
    \item \text{Optimization Problem}
    \item \text{QUBO}
    \item \text{Hamiltonian}
    \item \text{Ising Model}
    \item \text{Quantum Annealing}
    \item \text{Ground State}
    \item \text{Optimal Solution}
\end{enumerate}
This transformation is the fundamental idea behind quantum optimization algorithms and quantum annealing hardware.